\chapter{Conclusiones}
\label{cap:conclusiones}

\section{Respuesta a la pregunta de investigación}

La elección del método se sustenta en la estructura del cotizador y en el costo de
su implementación. La caracterización analítica obtenida puede resumirse así:

\begin{aporte}
En el modelo estudiado, la valoración es \textbf{afín a trozos}, con umbrales
candidatos organizados en tres racimos geométricos para la parametrización
considerada. La extensión a $M\ge0$ es constante en $[0,\SMLV]$, mientras que en
el dominio admisible $M\ge\SMLV$ hay una raíz única si y solo si
$\KA\ge\VP(\SMLV)$. El auxilio funerario puede destruir la convexidad. Las
mesadas de la muestra de renta vitalicia examinada se sitúan en un intervalo
afín; esa observación no describe la cartera completa de la compañía ni el
mercado colombiano.
\end{aporte}

La derivada lateral exacta permite aprovechar los tramos sin aproximar pendientes
por valoraciones auxiliares. Un paso de Newton obtiene la raíz cuando la clausura
del tramo seleccionado contiene también la solución. Esta propiedad no garantiza
una o dos iteraciones desde todo arranque ni identifica los pasos con el número
de tramos atravesados. Para el problema interior con un encaje inicial válido, la
variante semisuave con bisección periódica aporta una garantía global explícita
(\cref{thm:salvaguardado}), por lo que constituye una recomendación fundada para
el cotizador estudiado. Su costo y el de las otras alternativas deben leerse en
el \cref{cap:resultados}, con la contabilidad completa allí definida; la
recomendación no establece una superioridad universal.

\section{Verificación de las hipótesis de trabajo}

\begin{description}[leftmargin=1.1cm,style=nextline]
\item[H1 --- caracterización analítica.] La representación afín a trozos y el
  conjunto de hasta $\datoNQuiebresMax$ umbrales candidatos se obtienen de las
  operaciones de máximo y recorte (\cref{thm:quiebres}). Un umbral candidato
  puede no producir un salto efectivo si su peso es nulo.

\item[H2 --- terminación bajo condiciones precisas.] La inversión exacta de un
  tramo da la terminación en un paso si este contiene la raíz en su clausura.
  La cota por número de tramos requiere las condiciones adicionales del
  \cref{thm:terminacion-finita}; no es una igualdad general. Los experimentos
  ilustran estas condiciones, sin convertir el comportamiento de la muestra en
  una garantía para toda cotización.

\item[H3 --- contabilidad de costo.] Las iteraciones no representan el mismo
  trabajo en todos los métodos. La comparación contabiliza preparación,
  valoraciones auxiliares y llamadas efectivas al motor, y las complementa con
  tiempo de ejecución. Los cocientes de costo corresponden al protocolo y al
  entorno medidos.

\item[H4 --- derivada analítica conjunta.] Reutilizar los vectores de valoración
  evita las llamadas auxiliares de las diferencias finitas. El costo marginal de
  las derivadas no se considera nulo, y el ahorro efectivo se evalúa con el oráculo
  conjunto de valor y derivadas. No se deduce un porcentaje fijo del número
  nominal de evaluaciones por iteración.
\end{description}

\section{Contribuciones}

\subsection{Al análisis del problema}

El aporte central es integrar herramientas conocidas de valoración y análisis no
suave para estudiar este cotizador. La condición de capital mínimo distingue
falta de financiación, solución en el piso y solución interior. La expresión de
las derivadas laterales da acceso al $B$-subdiferencial y relaciona el error de
reserva con el de mesada. Las condiciones de terminación y el encaje separan la
rapidez local de la garantía global.

En retiro programado, el estado incluye el saldo y la mesada anterior porque los
ajustes dependen de ambos. La cota del \cref{thm:propagacion} incorpora ese
acoplamiento y permite escoger tolerancias a partir de un error admisible del
saldo final, bajo sus hipótesis. No implica una amplificación universal ni
permite comparar la conveniencia de las modalidades pensionales.

\subsection{A la práctica}

La implementación incorpora las derivadas laterales, un catálogo instrumentado de
\datoNMetodos{} métodos y un protocolo que distingue trabajo del motor,
precisión y tiempo. Los hallazgos de la \cref{sec:hallazgos-codigo} permiten
revisar el tratamiento del piso, la parada y los ajustes de retiro programado.
La aplicación interactiva facilita explorar las consecuencias de los parámetros;
su uso no sustituye una nota técnica ni la validación de un producto comercial.

\section{Limitaciones}

\begin{enumerate}[leftmargin=*]
  \item \textbf{Alcance del modelo.} Se analiza una implementación particular, no
        una representación exhaustiva de las obligaciones regulatorias ni de la
        tarificación comercial. Reserva, capital y prima requieren conciliación
        sobre fecha, componentes y base de valoración comunes.

  \item \textbf{Representatividad.} Las muestras proceden de una sola compañía y
        no permiten inferir propiedades de su cartera completa o del mercado.
        Los escenarios sintéticos amplían la cobertura numérica, sin aportar
        representatividad estadística de esas poblaciones.

  \item \textbf{Bases actuariales.} La procedencia de las series de mortalidad,
        los años base de mejora y los factores de aplicabilidad requieren
        documentación adicional. El ajuste descriptivo de Gompertz no valida
        esos niveles; la dependencia entre vidas tampoco se cuantifica. Cambiar
        pesos independientes de $M$ conserva la forma afín a trozos, pero las
        propiedades que requieren positividad deben comprobarse de nuevo.

  \item \textbf{Calibración y retiro programado.} La volatilidad, los recargos de
        mortalidad y el beneficiario hipotético no cuentan con calibración
        acreditada. El factor NI puede reproducir la referencia empleada para
        calibrarlo, lo que no constituye validación independiente. Deben
        conciliarse además las bases monetarias de rendimientos y pagos y la
        aproximación temporal de los retiros. Las trayectorias son ilustrativas.

  \item \textbf{Condiciones de la cota.} La propagación supone trayectorias con
        las mismas funciones de valoración, parámetros y rendimientos, y
        soluciones admisibles en un horizonte común. Cambiar referencias o
        recalibrar entre trayectorias exige términos adicionales. El cierre de
        una simulación por insuficiencia de capital no determina la extinción
        jurídica del derecho pensional.

  \item \textbf{Garantías y medición.} La terminación exacta se formula en
        aritmética exacta; el programa se detiene con tolerancias y redondeo.
        Los tiempos dependen del entorno y los órdenes empíricos requieren
        suficientes errores antes de la saturación numérica. La reproducción
        exige conservar versiones, entradas y configuración de cada ejecución.
\end{enumerate}

\section{Trabajo futuro}

\begin{enumerate}[leftmargin=*]
  \item \textbf{Conciliación actuarial.} Documentar fuentes, transformaciones,
        fechas y componentes de valoración; explicar las discrepancias antes de
        introducir multiplicadores y separar datos de calibración y evaluación.

  \item \textbf{Curvas de descuento.} Calibrar Smith--Wilson con observaciones y
        convenciones monetarias coherentes, contrastar alternativas y examinar
        su efecto en las reservas. Contar con más nodos no acredita por sí solo
        la idoneidad de la curva.

  \item \textbf{Dependencia y experiencia.} Estimar conjuntamente las bases de
        mortalidad y su posible dependencia con datos adecuados. Si la
        dependencia modifica solo pesos independientes de $M$, puede estudiarse
        su efecto sin perder la representación afín a trozos.

  \item \textbf{Proyección y comparación de modalidades.} Justificar la dinámica
        financiera y los ajustes, utilizar trayectorias comunes para comparar
        sensibilidades y definir criterios de riesgo y suficiencia antes de
        extraer conclusiones sobre modalidades pensionales.

  \item \textbf{Arranque en caliente.} Medir su utilidad y establecer condiciones
        para que la mesada anterior, expresada en la base monetaria pertinente,
        permanezca en un tramo que contenga la nueva raíz. La cercanía entre años
        no garantiza esa condición.

  \item \textbf{Extensiones del problema.} Estudiar la inversión respecto de la
        tasa y las bases de mortalidad estocásticas. La forma afín a trozos en la
        mesada no se traslada automáticamente a otras incógnitas; la teoría
        necesaria depende de si se busca la raíz de una esperanza o una raíz
        condicionada a cada escenario.
\end{enumerate}

\section{Reflexión final}

La \cref{fig:residuo-lineal} ilustra por qué una función cercana a una recta puede
requerir atención en sus quiebres y en el piso. Reconocer esa estructura permite
resolver el caso habitual con pocas operaciones y, a la vez, identificar los
casos donde faltan capital, hipótesis de convergencia o evidencia actuarial. El
valor del análisis reside en delimitar esas condiciones y documentar la precisión
obtenida, sin confundir exactitud numérica con validez del modelo.
